\documentclass[10pt,a4paper]{amsart}
\usepackage[margin=19mm]{geometry}
\usepackage{amsmath,amssymb,mathtools}
\usepackage[scr=rsfs,cal=euler]{mathalfa}
\usepackage{xfrac}
\usepackage[unicode,hidelinks,bookmarks=false]{hyperref}
\newcommand{\ie}{i.e.\ }
\newcommand{\cf}{cf.\ }
\newcommand{\resp}{respectively\ }
\newcommand{\Subsection}{\subsection}
\providecommand{\nobreakdash}{\nobreak\hbox{-}}
\setlength{\emergencystretch}{2em}
\multlinegap0pt
\usepackage{bm}
\usepackage{bbm}
\usepackage{dsfont}
\usepackage{xspace}
\usepackage{cancel}
\usepackage{mathtools}
\usepackage{amsthm}
\makeatletter
\@ifundefined{theorem}{\newtheorem{theorem}{Theorem}}{}
\@ifundefined{lemma}{\newtheorem{lemma}[theorem]{Lemma}}{}
\makeatother
\title[Discrete Space-Time Wave Kernels on Regular Trees]{Discrete Space-Time Wave Kernels on Regular Trees}
\author{Amar Ba\v{s}i\'c, Lejla Smajlovi\'c, Zenan \v{S}abanac}
\date{Source: arXiv:2606.27083v1 (CC BY 4.0)}
\begin{document}
\maketitle

\noindent\emph{Source and licence.} Discrete Space-Time Wave Kernels on Regular Trees. Version arXiv:2606.27083v1 (CC BY 4.0), distributed under the Creative Commons Attribution 4.0 International licence, as recorded in the arXiv licence metadata for this exact version and shown on the abstract page https://arxiv.org/abs/2606.27083v1. Full rights notice: licenses/arxiv-2606.27083-CC-BY-4.0.txt.\par\smallskip

\noindent\textit{Editorial scope.} Counted unit: the Lemma whose source label is \texttt{lem: kernels on tree} in the article (source0.tex lines 389--399, source label \texttt{lem: kernels on tree}), delivered together with the article's own complete original proof of it (source0.tex lines 400--434, one \texttt{proof} environment of 35 lines). No mathematics is written, repaired or re-derived: statement and proof are the article's own text line for line. Required context retained uncounted: wave\_eq\_disc (lines 145--150); cond1 (lines 168--174); cond on V (lines 176--182). Every definition, display and label of those lines is retained unchanged. Omitted from this package and not redistributed: 1--144, 151--167, 175--175, 183--388, 435--1394. Nothing inside a retained excerpt is omitted or altered: every retained body is delivered exactly as the article's lines are written, so no omission receipt of any kind accompanies this package. Citations: the retained excerpts carry no citation command at all, so no bibliography entry of the article is transcribed and no reference is redistributed. Reference adaptation: none was needed and none was made; every reference inside the delivered unit resolves inside it, so no unresolved reference or citation is produced. Printed slips kept literal: no printed word, symbol, hypothesis or number of the counted statement or of its proof was altered. Layout and class: the article's own class is \texttt{amsart} with packages including inputenc, fontenc, amsfonts, amsmath, eucal, amssymb, amstext, amsthm, mathrsfs, centernot, enumitem, graphicx, color, caption; the wrapper is the lane's pinned \texttt{amsart} editorial wrapper at 10pt on A4, which restates the theorem-like environments the retained text needs and adds the article's own macro definitions that the retained text needs (none), with extra wrapper packages (bm, bbm, dsfont, xspace, cancel, mathtools). The delivered numbering is the unit's own. Assets: no retained span contains a figure environment, a graphics-inclusion command, a PDF-page inclusion, a table or a TikZ picture; no image file is copied into this package, the package declares no asset dependency and no image file is redistributed with it. Licence: version arXiv:2606.27083v1 of this article is distributed under the Creative Commons Attribution 4.0 International licence, as recorded in the arXiv licence metadata for this exact version and shown on the abstract page https://arxiv.org/abs/2606.27083v1. A full rights notice is delivered at licenses/arxiv-2606.27083-CC-BY-4.0.txt. Characters: every retained excerpt is pure ASCII, so the delivered engine, XeLaTeX with its default Latin Modern text font, reports no missing character.
\begin{equation}\label{wave_eq_disc}
\Delta_{q+1}^{a,b}\, W(x_0, x; t)
\;+\;
\partial_{t,\mathbb{T}}^2 W(x_0, x; t)
= 0
\end{equation}

\begin{equation}\label{cond1}
W_{q+1}^{a,b}(x_0, x; 0) =\delta_{x_0}(x):=
\begin{cases}
1, & \text{if } x = x_0, \\
0, & \text{if } x \neq x_0,
\end{cases} \quad \text{and}\quad \partial_t W_{q+1}^{a,b}(x_0, x; 0) = 0,\,\, \forall x\in V(T_{q+1}),
\end{equation}

\begin{equation}\label{cond on V}
V_{q+1}^{a,b}(x_0, x; 0) = 0 \quad \text{and}\quad  \partial_t  V_{q+1}^{a,b}(x_0, x; 0) =
\begin{cases}
1, & \text{if } x = x_0, \\
0, & \text{if } x \neq x_0,
\end{cases} \,\, \forall x\in V(T_{q+1}).
\end{equation}

\begin{lemma} \label{lem: kernels on tree}
For any real parameters $a,\, b$, with $b\neq 0$, the functions
\begin{equation}\label{wave w as conv}
W_{q+1}^{a,b}(|x|;t) = \sum_{k=0}^{\infty} (-1)^k h_{2k}(t,0)\,(q+1)^k \left(\left(b-\frac{a}{q+1}\right) \delta_0 -b\nu \right)^{(*k)}(x),\,\, x\in T_{q+1},
\end{equation}
and
\begin{equation} \label{wave V as conv}
V_{q+1}^{a,b}(|x|;t) = \sum_{k=0}^{\infty} (-1)^k h_{2k+1}(t,0)\,(q+1)^k \left(\left(b-\frac{a}{q+1}\right) \delta_0 -b\nu \right)^{(*k)}(x),\,\, x\in T_{q+1},
\end{equation}
satisfy equation \eqref{wave_eq_disc} and the initial conditions \eqref{cond1} and \eqref{cond on V}, respectively.
\end{lemma}

\begin{proof}
First, note that the coefficients $h_{2k}(t,0)=\binom{t}{2k}$ satisfy the binomial difference identity
$$
\partial_t^2 h_{2k}(t,0)=h_{2k-2}(t,0)\quad(k\ge1), \qquad \partial_t^2 h_0(t,0)=0.
$$
For brevity, set $B=b-\frac{a}{q+1}$. Since the sums defining $W_{q+1}^{a,b}(|x|;t)$ and $V_{q+1}^{a,b}(|x|;t)$ are finite for fixed $t$, we may interchange the summation with the action of the operators $\partial_t$ and $\Delta_{q+1}$, which yields 
\begin{equation*}
\begin{split}
\partial_t^2 W_{q+1}^{a,b}(|x|;t)
&= \sum_{k=1}^{\infty}(-1)^k h_{2k-2}(t,0)(q+1)^k \left(B\delta_0 -b\nu \right)^{(*k)}(x)\\
&=-\sum_{k=0}^{\infty}(-1)^k h_{2k}(t,0)(q+1)^{k+1} \left( B \delta_0 -b\nu \right)^{(\ast(k+1))}(x)\\
&=
-(q+1)\left(B \delta_0 -b\nu \right)\ast \sum_{k=0}^{\infty} (-1)^k h_{2k}(t,0)\,(q+1)^k \left(B\delta_0 -b\nu \right)^{(*k)}(x) \label{eq. wave eq. prop}\\
&=
-\Delta_{q+1}^{a,b}W_{q+1}^{a,b}(|x|;t).
\end{split}
\end{equation*}
Hence, $W_{q+1}^{a,b}(|x|;t)$ satisfies \eqref{wave_eq_disc}.

\medskip
At $t=0$, we have $h_0(0,0)=1$ and $h_{\ell}(0,0)=0$ for all $\ell\ge1$, so
$$
W_{q+1}^{a,b}(|x|;0)=\left( B\delta_0 - b\nu \right)^{(*0)}(x)= \delta_{x=x_0}(x)
$$
and
$$
\partial_t W_{q+1}^{a,b}(|x|;0)=\sum_{k=1}^{\infty} (-1)^k h_{2k-1}(0,0)\,(q+1)^k \left( B\delta_0 - b\nu \right)^{(*k)}(x)\equiv 0.
$$

A similar argument shows that $V_{q+1}^{a,b}(|x|;t)$ satisfies \eqref{wave_eq_disc}. To show that the initial conditions \eqref{cond on V} are satisfied, we use that $h_{2k+1}(0,0)=0$ for all $k\geq 0$ to see that $V_{q+1}^{a,b}(|x|;0)=0$, while delta differentiation of $V_{q+1}^{a,b}$ with respect to $t$ yields that
$$
\partial_t V_{q+1}^{a,b}(|x|;0)=W_{q+1}^{a,b}(|x|;0)=\delta_0(x),
$$
which proves \eqref{cond on V}.
\end{proof}

\end{document}
